% tex/impl/commands/rules.tex
\ProvidesFile{rules.tex}[2026/01/04 ND rule macros]

% Erwartung:
% - hyperref ist bereits geladen (in der Preambel, nicht hier)
% - Labels heißen: rule:<NAME>  (z.B. rule:AI, rule:RE, ...)

% ------------------------------------------------------------
% Zentrales Helfer-Makro: erzeugt den klickbaren Regelnamen
% ------------------------------------------------------------
\providecommand{\RuleRef}[2]{%
  \hyperref[rule:#1]{\ensuremath{#2}}%
}

% Link only the rule name, not its justification argument.  Arguments often
% contain theorem or rule references themselves; keeping them outside the
% outer link avoids invalid nested PDF links while preserving every link.
\providecommand{\RuleRefWithArg}[3]{%
  \ensuremath{\hyperref[rule:#1]{#2}(#3)}%
}

% ------------------------------------------------------------
% Annahme
% ------------------------------------------------------------
\providecommand{\rA}{\RuleRef{A}{A}}

% ------------------------------------------------------------
% Konjunktion ∧
% ------------------------------------------------------------
\providecommand{\rAI}[1]{\RuleRefWithArg{AI}{\land I}{#1}}
\providecommand{\rAEa}[1]{\RuleRefWithArg{AE1}{\land E1}{#1}}
\providecommand{\rAEb}[1]{\RuleRefWithArg{AE2}{\land E2}{#1}}
\providecommand{\rAEn}[1]{\RuleRefWithArg{AEn}{\land E}{#1}}

% ------------------------------------------------------------
% Disjunktion ∨
% ------------------------------------------------------------
\providecommand{\rOIa}[1]{\RuleRefWithArg{OI1}{\lor I1}{#1}}
\providecommand{\rOIb}[1]{\RuleRefWithArg{OI2}{\lor I2}{#1}}
\providecommand{\rOE}[1]{\RuleRefWithArg{OE}{\lor E}{#1}}

% ------------------------------------------------------------
% Implikation →
% ------------------------------------------------------------
\providecommand{\rRI}[1]{\RuleRefWithArg{RI}{\rightarrow I}{#1}}
\providecommand{\rRE}[1]{\RuleRefWithArg{RE}{\rightarrow E}{#1}}

% ------------------------------------------------------------
% Äquivalenz ↔
% ------------------------------------------------------------
\providecommand{\rLRI}[1]{\RuleRefWithArg{LRI}{\leftrightarrow I}{#1}}
\providecommand{\rLREa}[1]{\RuleRefWithArg{LRE1}{\leftrightarrow E1}{#1}}
\providecommand{\rLREb}[1]{\RuleRefWithArg{LRE2}{\leftrightarrow E2}{#1}}
\providecommand{\rLRS}[1]{\RuleRefWithArg{LRSubst}{\leftrightarrow S}{#1}}

% ------------------------------------------------------------
% Allquantor ∀
% ------------------------------------------------------------
\providecommand{\rUI}[1]{\RuleRefWithArg{UI}{\forall I}{#1}}
\providecommand{\rUE}[1]{\RuleRefWithArg{UE}{\forall E}{#1}}

% ------------------------------------------------------------
% Existenzquantor ∃
% ------------------------------------------------------------
\providecommand{\rEI}[1]{\RuleRefWithArg{EI}{\exists I}{#1}}
\providecommand{\rEE}[1]{\RuleRefWithArg{EE}{\exists E}{#1}}

% ------------------------------------------------------------
% Eindeutige Existenz ∃!
% (Du nutzt aktuell \UEI/\UEE ohne führendes r; \rUEI bleibt als Alias.)
% ------------------------------------------------------------
\providecommand{\UEI}[1]{\RuleRefWithArg{UEI}{\exists! I}{#1}}
\providecommand{\UEE}[1]{\RuleRefWithArg{UEE}{\exists! E}{#1}}

\providecommand{\rUEI}[1]{\UEI{#1}}

% ------------------------------------------------------------
% Gleichheit =
% ------------------------------------------------------------
\providecommand{\rII}{\RuleRef{II}{=I}}
\providecommand{\rIE}[1]{\RuleRefWithArg{IE}{=E}{#1}}

% ------------------------------------------------------------
% Negation ¬ und ⊥
% ------------------------------------------------------------
\providecommand{\rBI}[1]{\RuleRefWithArg{BI}{\bot I}{#1}}
\providecommand{\rCI}[1]{\RuleRefWithArg{CI}{\neg I}{#1}}
\providecommand{\rCE}[1]{\RuleRefWithArg{CE}{\neg E}{#1}}

% ------------------------------------------------------------
% Kettennotation / Transitivitätsschritt
% ------------------------------------------------------------
\providecommand{\rChain}[1]{\RuleRefWithArg{Chain}{\mathsf{Tr.}}{#1}}

% ------------------------------------------------------------
% Induktion (falls du das als "Regel" klickbar willst)
% ------------------------------------------------------------
\providecommand{\rInduktion}[1]{\RuleRefWithArg{Induktion}{\mathrm{Induktion}}{#1}}
